
La comparación reduce el problema a controlar las reparaciones a lo largo del
proceso. Mientras quede una fracción fija de piezas, los grados y los normalizadores
se concentran alrededor de sus trayectorias deterministas (Sección~\ref{sec:regularidad}).
Un \emph{patrón} fija qué se exige a las casillas cuya reina se retira, en una reparación o
en su operación inversa, y a las diagonales de las casillas que se ocupan: que cada casilla
tenga ya su reina o siga disponible, y que cada diagonal esté libre de toda pieza o, al
menos, de reinas. Para la cuenta directa seguimos $32$ patrones
por cada casilla candidata de un cruce, y para la inversa $16$ por cada par de casillas.
En cada familia, el cambio esperado de un patrón se expresa mediante los patrones de la
misma familia, con un error que se acota en la prueba. Resolviendo hacia atrás el sistema
lineal correspondiente se obtiene una combinación lineal de patrones, el \emph{compensador},
que coincide con el contador en el instante elegido y cuya deriva es pequeña; una
desigualdad de martingalas da entonces la concentración hasta un corte
temprano, cuando quedan unas $\delta n$ reinas por colocar. Después del corte, cada reina
cambia el contador directo de cada cruce en a lo sumo cinco unidades y el contador
inverso en $O(n)$. Al dividir por $n$ y por $n^2$, respectivamente, el efecto
de todo el tramo final es $O(\delta)$.

Una vez controlada la pérdida tras el corte, la desigualdad de entropía binaria
transfiere la concentración a la ley uniforme de cada capa. Para la cota inferior
basta la cuenta directa, pues con probabilidad tendiente a uno cada cruce vacío
admite al menos $a_0n/2$ reparaciones. Para hallar el límite hay que tener en
cuenta la posición del cruce. En $(r,c)$, la cuenta directa se concentra alrededor
de $n\,a(\bar r,\bar c)$. Al elegir el peso $w=1/a$ obtenemos
$D_{1/a}/(h^2n)\to1$. La media del contador inverso dividido por $n^2$, con ese mismo
peso, tiende a $\int\Gamma/a$,
donde $\Gamma(r,c)$ es la densidad límite de las operaciones inversas que crean el cruce $(r,c)$
(Lema~\ref{lem:nucleo}), y la identidad $\Gamma=b_0\,r_Rr_C\,a$, consecuencia de la forma producto del
perfil, da $\int\Gamma/a=b_0$. Por tanto $f_h\to1/b_0=F(0)$.
\ifverificacionLean
La Sección~\ref{sec:verificacion} describe la formalización en Lean de los enunciados
principales y precisa qué comprueba.
\fi

\begin{figure}[t]
\centering
\begin{tikzpicture}[
 caja/.style={draw,rounded corners=2pt,align=center,font=\footnotesize,inner sep=3pt,
 minimum height=9mm,text width=29mm},
 flecha/.style={-{Stealth[length=2mm]},thick}]
\node[caja] (perfil) at (0,0) {Perfil variacional\\ y primer orden (§\ref{sec:preliminares})};
\node[caja] (ref) at (3.6,0) {Comparación\\ entrópica (§\ref{sec:referencia})};
\node[caja] (din) at (7.2,0) {Regularidad\\ dinámica (§\ref{sec:regularidad})};
\node[caja] (abs) at (3.6,-1.7) {32 patrones directos\\ y corte (§\ref{sec:absorbedores})};
\node[caja] (inv) at (7.2,-1.7) {16 patrones\\ inversos (§\ref{sec:inversa})};
\node[caja] (dc) at (0,-3.4) {Doble cuenta\\ entre capas (§\ref{sec:doblecuenta})};
\node[caja] (pos) at (3.6,-3.4) {Teorema~\ref{thm:positiva}\\ (§\ref{sec:positiva})};
\node[caja] (lim) at (7.2,-3.4) {Teoremas~\ref{thm:limite} y~\ref{thm:estructura}\\ (§\ref{sec:limite})};
\node[caja] (cor) at (10.8,-3.4) {Corolarios~\ref{cor:cociente}\\ y~\ref{cor:capas} (§\ref{sec:corolarios})};
\draw[flecha] (perfil) -- (ref);
\draw[flecha] (ref) -- (din);
\draw[flecha] (din) -- (abs);
\draw[flecha] (din) -- (inv);
\draw[flecha] (abs) -- (pos);
\draw[flecha] (abs) -- (lim);
\draw[flecha] (inv) -- (lim);
\draw[flecha] (dc) -- (pos);
\draw[flecha] (dc.south) to[out=-25,in=-155] (lim.south);
\draw[flecha] (lim) -- (cor);
\end{tikzpicture}
\caption{Estructura lógica de la demostración. La doble cuenta es puramente combinatoria; la
comparación entrópica de la Sección~\ref{sec:referencia} interviene en todas las flechas que
llegan a resultados bajo $\mu_h$.}
\label{fig:esquema}
\end{figure}

\subsection{Relación con trabajos previos}

Simkin~\cite{Simkin} caracteriza el término exponencial de $Q(n)$ mediante una optimización de
entropía y establece grandes desviaciones para el perfil de las soluciones completas; antes,
Luria~\cite{Luria} había obtenido cotas superiores con el método de entropía. La cota inferior
de primer orden que usamos procede del resultado de Simkin (Teorema~\ref{thm:simkin}). Para
probarla, Simkin fija una \emph{queenon} escalonada, es decir, una queenon
(Sección~\ref{sec:preliminares}) de densidad constante en cada bloque de una cuadrícula; en cada
paso elige un bloque con la probabilidad que esa medida le asigna, coloca en él una
reina al azar entre las casillas disponibles y, al final, completa la colocación con
absorbedores~\cite[Sección~5]{Simkin}. Luria y Simkin~\cite{LS} construyen casi-soluciones
mediante un proceso voraz aleatorio y las completan con absorbedores para obtener una cota
inferior de conteo, y Bowtell y Keevash~\cite{BK} estudian el conteo toroidal y sus
consecuencias para el tablero ordinario. El proceso de referencia y los absorbedores de este
trabajo siguen ese espíritu, pero el éxito bajo un proceso de construcción no identifica por sí
solo la ley uniforme sobre las colocaciones con $h$ huecos: esa identificación es el papel de
la comparación entrópica. Comparar la ley uniforme con la del resultado de un proceso aleatorio
es la herramienta principal del trabajo de Kwan~\cite{Kwan} sobre sistemas triples de Steiner.
A su vez, la doble cuenta de reparaciones entre capas consecutivas es una forma del método de
conmutaciones (\emph{switchings})~\cite{McKayWormald,HMcK}.

Completar una colocación parcial dada es un problema que se remonta a 1850 y es
NP-completo~\cite{GJN}. Glock, Munhá Correia y
Sudakov~\cite{GMS} prueban que toda colocación legal de a lo sumo $n/60$ reinas puede
completarse y construyen colocaciones no completables de tamaño cercano a $n/4$;
Nielsen~\cite{Nielsen} mejora la cota superior del umbral extremal a $0{,}216\,n$ para $n$
suficientemente grande. Esos
enunciados son de peor caso; aquí se estudian probabilidades bajo la ley uniforme de
colocaciones casi completas.

Hasta donde sabemos, los
Teoremas~\ref{thm:limite} y~\ref{thm:estructura} y el Corolario~\ref{cor:cociente} son nuevos.
No se reclama novedad para las herramientas generales: regla de la cadena de la entropía
relativa, comparación con procesos aleatorios~\cite{Kwan}, dobles cuentas y
conmutaciones~\cite{HMcK}, método de ecuaciones diferenciales~\cite{Wormald} y desigualdades de
concentración~\cite{AlonSpencer}.
